% Part: counterfactuals
% Chapter: minimal-change-semantics
% Section: introduction

\documentclass[../../../include/open-logic-section]{subfiles}

\begin{document}

\olfileid{cnt}{min}{int}

\olsection{Pambuka}

Stalnaker lan Lewis ngajokake panjelasan tumrap kondisional
kontrafaktual kaya ``Saupama penthol korek digesek, mesthi bakal
murub.'' Panjelasane minangka usulan kanggo mangerteni kanthi
bener syarat bebener ukara kaya mangkono. Gagasan kang ndhasari
loro usulan iku yaiku: kanggo nemtokake apa kondisional
kontrafaktual bener, kita kudu nimbang donya-donya kang bisa ana
lan kang bedane saka donya aktual paling sithik supaya
antesedene bener. Yen konsekuene bener ing donya-donya iku,
mula kontrafaktuale bener. Upamane, aku nyekel penthol korek
lan bungkus korek awujud buku. Ing donya aktual, aku mung
ndeleng lan mikir apa kang bakal kelakon saupama aku nggesek
penthol korek. Owah-owahan minimal saka donya aktual menyang
donya ing ngendi aku nggesek korek yaiku aku mutusake tumindak
lan nggesek korek. Iki minimal amarga ora ana owah-owahan
liyane: aku ora uga mlumpat, nggesek korek ora uga ngobong
rambutku, drijiku ora dumadakan kelangan kabeh kakuwatane,
aku ora ing wektu kang padha disiram banyu nalika disergap
nganggo SuperSoaker, lan sapiturute. Ing kemungkinan alternatif
iku, penthol korek murub. Mula bener yen saupama aku nggesek
penthol korek, penthol iku mesthi bakal murub.

Panjelasan intuitif iki bisa digandhengake karo semantik formal
kanggo logika kontrafaktual. Lewis ngenalake simbol ``$\boxright$''
kanggo kontrafaktual, dene Stalnaker nggunakake simbol~``$>$''.
Kita bakal nggunakake~$\cif$ lan nambahake minangka konektif
biner ing logika proposisional. Dadi, saliyane !!{formula}s
awujud $!A \lif !B$, kita uga nduweni !!{formula}s awujud~$!A \cif !B$.
Semantik formal iki, kaya semantik relasional logika modal,
adhedhasar model kang ngevaluasi !!{formula}s ing donya-donya,
lan syarat pemenuhan kang nemtokake $\mSat{M}{!A \cif !B}[w]$
diandharake nganggo $\mSat{M}{!A}[w']$ lan $\mSat{M}{!B}[w']$
kanggo sawenehing donya (liyane)~$w'$. Donya $w'$ kang endi?
Kanthi intuitif, donya kang paling cedhak karo~$w$ lan ing
kono $\mSat{M}{!A}[w']$ lumaku. Iki mbutuhake relasi
``cedhak'' uga dilebokake ing model.

Lewis ngenalake cara kang menehi pangerten kanggo nggambarake
kaanan kontrafaktual. Saben donya kang bisa ana ana ing pusat
himpunan bola-bola kang susun ing njero siji lan sijine lan
ngemot donya liyane---kita nggambar bola-bola iki minangka
bunderan kanthi pusat kang padha. Donya-donya ing lapisan
antarane rong bola padha cedhake karo donya pusat; donya ing
bola njero luwih cedhak, lan donya ing bola kang ngubengi
luwih adoh.
\begin{center}
\begin{tikzpicture}[scale=.7]\small
  \spheresystem{5}
  \draw (0,0) node {$w$};
  \propositionintersect{0}{4}{40}{3.5}
  {\tikzset{proposition/.append={smooth,tension=1.4}}
  \proposition[shift={(1.6,-1)}]{90}{1}{30}{4}}
  \path (1.5,3) node[below] {$\formula{B}$};
  \path (3,0) node {$\formula{A}$};
\end{tikzpicture}
\end{center}
Ing gambar iki, donya-donya $!A$ kang paling cedhak yaiku
donya-donya $w'$ kang nyukupi~$!A$ lan ana ing bola paling
cilik kang ngemot donya kaya mangkono ing saubenge donya
pusat~$w$ (wewengkon klawu). Kanthi intuitif, $!A \cif !B$
dipenuhi ing~$w$ yen $!B$ bener ing kabeh donya $!A$ kang
paling cedhak.

\end{document}
